\section{Limitations and Broader Impact}

The controlled evidence is limited to approximately 98.2M parameters, one tokenizer, one FineWeb-Edu shard, 100M training tokens, and three seeds. The paired ordering is consistent, but its 95\% interval includes zero. It does not establish statistical significance, scaling behavior, downstream-task benefit, or universal superiority to GQA.

The comparison changes a bundle: the R/W/V parameterization and per-head R/W RMS normalization. The six-run evidence does not isolate normalization factorially. Exploratory single-seed lexical-write and no-normalization observations are excluded from the controlled table because they are not three-seed matched comparisons. Additional matched controls are required before assigning causal credit to one component.

Training throughput is implementation-, software-, and H200-specific. Eager execution was used because the tested compiler stack produced non-finite gradients in a shared lexical-table path. The reported 2.86\% penalty is not a production inference estimate. W/R/V retains W/V tensors whose size grows with context, like a grouped KV cache. Long-context prefill, time to first token, serving latency distributions, memory bandwidth, and CUDA Graph replay remain outside the claims.

FineWeb-Edu inherits limitations of large web corpora, including uncertain provenance, representational imbalance, and potentially harmful content. The architecture does not remove these risks. Efficiency improvements can also lower the cost of deploying low-quality or harmful models. The archived dataset shard is included for exact reproducibility and should be handled under its source license and governance requirements.

Finally, the work is an architecture-path study rather than a claim of a novel softmax operator. R and W occupy the technical query/key boundary of SDPA. The value of the terminology is the explicit data-flow decomposition and executable comparison, not a claim that renaming Q/K/V alone is an algorithmic contribution.